\documentclass[10pt,a4paper]{amsart}
\usepackage[margin=19mm]{geometry}
\usepackage{amsmath,amssymb,mathtools}
\usepackage[scr=rsfs,cal=euler]{mathalfa}
\usepackage{xfrac}
\usepackage[unicode,hidelinks,bookmarks=false]{hyperref}
\newcommand{\ie}{i.e.\ }
\newcommand{\cf}{cf.\ }
\newcommand{\resp}{respectively\ }
\newcommand{\Subsection}{\subsection}
\providecommand{\nobreakdash}{\nobreak\hbox{-}}
\setlength{\emergencystretch}{2em}
\multlinegap0pt
\usepackage{amssymb}
\newtheorem{lemma}{Lemma}
\title{Norm rigidity and equality cases for the Dyn--Farkhi inequality}
\author{Mark Meyer}
\date{Source: arXiv:2608.13738v1 (CC BY 4.0)}
\begin{document}
\maketitle

\noindent\emph{Source and licence.} Norm rigidity and equality cases for the Dyn--Farkhi inequality. Version arXiv:2608.13738v1 (CC BY 4.0), distributed by arXiv under the Creative Commons Attribution 4.0 International licence, http://creativecommons.org/licenses/by/4.0/, as recorded in the arXiv licence metadata for this exact version and shown on the abstract page https://arxiv.org/abs/2608.13738v1. Full licence notice: \texttt{licenses/arxiv-2608.13738-CC-BY-4.0.txt}.\par\smallskip

\noindent\textit{Editorial scope.} Counted unit: the article's lemma lemma:one\_dimensional\_one\_non\_convex (source0.tex lines 192--198). The article's complete original proof of it is delivered with it (source0.tex lines 200--239, one \texttt{proof} environment of 40 lines). No mathematics is written, repaired or re-derived: the statement and the proof are the article's own lines, in the article's own notation, and no retained line is substituted or re-flowed.  No further source-local context is needed: every cross-reference of every retained line resolves inside the two delivered spans.  Nothing was omitted from any retained span, so the package carries no omission record.  No retained line contains a citation, so the wrapper carries no bibliography and no bibliography entry.  No retained line contains a figure environment, an image-inclusion command or drawing code, so no image is delivered and the package declares no asset dependency.  Editorial formatting only, in addition: an AMS \texttt{amsart} preamble that restates the article's own declarations and macros used by the retained lines, namely the environments \texttt{lemma} on separate counters, together with the standard packages those lines need.  The retained environments print in this document as Lemma 1 in order of appearance, whereas the article prints the same items with the numbering of its own full document.  The complete native source of the article as distributed in the arXiv e-print for this exact version is bundled unchanged as \texttt{sources/207608/source0.tex}.
\begin{lemma}\label{lemma:one_dimensional_one_non_convex}
    Let $A,B\subset\mathbb{R}^2$ be compact one-dimensional sets such that one of them is nonconvex. Then
    \begin{align*}
        d(A+B)^2=d(A)^2+d(B)^2
    \end{align*}
    if and only if $\textup{aff}(A)$ and $\textup{aff}(B)$ are orthogonal lines. 
\end{lemma}

\begin{proof}
    By translating $A$ and $B$ independently, we may assume that $0\in A\cap B$, and write $\textup{aff}(A)=\textup{span}(u)$ and $\textup{aff}(B)=\textup{span}(v)$ for unit vectors $u,v$.

    If $u\bot v$, then choose $ru\in \textup{conv}(A)$ and $sv\in \textup{conv}(B)$ that attain the supremum in the definition of $d(A)$ and $d(B)$. Then
    \begin{align*}
        d(ru+sv,A+B)^2=\inf_{r_1u\in A,s_1v\in B}(|r-r_1|^2+|s-s_1|^2)=d(A)^2+d(B)^2.
    \end{align*}
    Equality follows immediately.

    Next, assume that $c:=\langle u,v\rangle \neq 0$. We will show that the distance from any $x\in \textup{conv}(A+B)$ to $A+B$ is strictly smaller than $d(A)^2+d(B)^2$. First, suppose that both sets are nonconvex, and write $x=ru+sv$. If $x\in \textup{conv}(A)+B$ or if $x\in A+\textup{conv}(B)$, then
    \begin{align*}
        d(x,A+B)^2\leq \max\{d(A)^2,d(B)^2\}<d(A)^2+d(B)^2.
    \end{align*}
    Otherwise, choose
    \begin{align*}
        r_1<r<r_2,\qquad r_iu\in A,\qquad r_2-r_1\leq 2d(A),
    \end{align*}
    and similarly, choose $s_1< s< s_2$, $s_iv\in B$, $s_2-s_1\leq 2d(B)$. Let $s_j$ be the closest to $s$, and set
    \begin{align*}
        t:=r-r_1,\qquad \gamma:=r_2-r_1,\qquad q:=s-s_j.
    \end{align*}
    Then $0<t<\gamma \leq 2d(A)$ and $0<|q|\leq d(B)$. We claim that the squared distance from $x$ to at least one of the points $r_1u+s_jv$ and $r_2u+s_jv$ is strictly less than $d(A)^2+q^2$. If not, then
    \begin{align*}
        t^2+2tqc\geq d(A)^2,\qquad (\gamma-t)^2-2(\gamma-t)qc\geq d(A)^2.
    \end{align*}
    Divide the left by $t$, the right by $\gamma-t$, and then add to get
    \begin{align*}
        \gamma\geq d(A)^2\left(\frac{1}{t}+\frac{1}{\gamma-t}\right)\geq\frac{4d(A)^2}{\gamma}.
    \end{align*}
    Then $\gamma\geq 2d(A)$ and as a result the above inequalities are all equality. This forces $\gamma=2d(A)$, $t=d(A)$, and $qc=0$, which is a contradiction. Therefore,
    \begin{align*}
        d(x,A+B)^2<d(A)^2+q^2\leq d(A)^2+d(B)^2.
    \end{align*}

    To complete the proof, we consider the case where $A$ is convex and $B$ is nonconvex, and the other case follows by a symmetry argument. Then $x=ru+sv$, with $ru\in A$. If $sv\in B$, the distance from $x$ to $A+B$ is zero. Otherwise, choose $s_1,s_2$ as above. The closest point $s_jv$ has distance to $B$ strictly less than $d(B)$ unless $s$ is the midpoint between $s_1,s_2$ and $s_2-s_1=2d(B)$. In this case, choose another point $ru+\varepsilon \delta u$ in $A$ in direction $\delta u$, $\delta\in \{-1,1\}$ at small distance $\varepsilon$ from $ru$, and choose $s_jv$ so that $\delta(s-s_j)c>0$. If $\varepsilon <2|(s-s_j)c|$, then
    \begin{align*}
        \|ru+sv-((ru+\varepsilon \delta u )+s_jv)\|_2^2=\varepsilon^2+d(B)^2-2\varepsilon|(s-s_j)c|<d(B)^2.
    \end{align*}
    By the continuity of the norm $\|\cdot \|_2$ and the compactness of $A+B$, the distance $d(A+B)$ is attained for some $x\in \textup{conv}(A+B)$. Therefore, the pointwise strict inequality that we proved above is actually a strict inequality for $d(A+B)$, so equality is impossible.
\end{proof}

\end{document}
